\chapter{Troubleshooting}
\label{app:troubleshooting}

\noindent
Symptom, cause, fix. Organised by what you observe rather than by what is wrong,
because when you need this page you do not yet know what is wrong.

\chapterstub{
\textbf{Contents.} Grouped by symptom. (C.1) The service becomes unresponsive
under load. (C.2) \code{Task exception was never retrieved}, and its silent
cousin, the garbage-collected task. (C.3) The process will not shut down.
(C.4) \code{coroutine has no attribute}, and mixing sync and async paths.
(C.5) \api{InvalidUpdateError} from parallel writes. (C.6) State not persisting
between turns --- almost always a missing \code{thread\_id}. (C.7) The
conversation vanished after a restart --- almost always the in-memory saver.
(C.8) \api{GraphRecursionError} and the agent that will not stop. (C.9) The same
side effect happened twice. (C.10) The model calls the wrong tool. (C.11) The
model stopped calling tools at all as the tool count grew. (C.12) Context
overflows, or quality degrades before it overflows. (C.13) Costs rose without a
traffic change. (C.14) Streaming works locally and buffers in production.
(C.15) Rate limits and 429s. (C.16) A resumed interrupt re-ran work.
(C.17) Evaluation scores moved and nothing was deployed.

\textbf{Convention.} Every entry states the observable symptom first, then the
cause, then the fix, then a cross-reference to the chapter that explains why.
Where a symptom has several causes, they are ordered by how often each turns out
to be the one.
}
